% TOP_PROBLEM: {"id": 2177, "title": "Erdős Problem #560", "category_id": 3, "category": {"id": 3, "name": "graph_theory", "display_name": "Graph Theory", "description": "Problems involving graphs, networks, and their properties.", "slug": "graph-theory", "order_index": 3, "created_at": "2026-07-31T15:26:25.671Z"}, "status": "open"}
% FIRST_POSTED: null
% ENTRY_KIND: statement_only
% Imported from: unsolvedmath_additional_500.tex; UnsolvedMath ID 2177
% !TeX program = lualatex
% Compile twice: lualatex 2177.tex
% Self-contained: no external bibliography, images, or \input files.
% Numeric section labels and visible numbers are original UnsolvedMath IDs.
\documentclass[11pt,a4paper]{article}
\usepackage[margin=24mm,headheight=14pt]{geometry}
\usepackage{fontspec}
\setmainfont{Latin Modern Roman}
\setsansfont{Latin Modern Sans}
\setmonofont{Latin Modern Mono}
\usepackage{amsmath,amssymb,mathtools}
\usepackage{unicode-math}
\setmathfont{Latin Modern Math}
\usepackage{microtype}
\usepackage{enumitem,xurl,needspace}
\usepackage[dvipsnames]{xcolor}
\usepackage{hyperref}
\hypersetup{unicode=true,colorlinks=true,linkcolor=MidnightBlue,urlcolor=MidnightBlue,
 pdftitle={UnsolvedMath Problem 2177},pdfauthor={Source-annotated compilation},bookmarksnumbered=true}
\usepackage{bookmark,tocloft,titlesec,fancyhdr}
\setlength{\cftsecnumwidth}{4.2em}
\cftsetpnumwidth{2.5em}
\cftsetrmarg{4em}
\titleformat{\section}{\Large\bfseries\raggedright}{\thesection}{0.7em}{}
\pagestyle{fancy}\fancyhf{}
\fancyhead[L]{\small\sffamily Additional UnsolvedMath statements}
\fancyhead[R]{\small Original catalogue IDs}
\fancyfoot[C]{\thepage}
\setlength{\parindent}{0pt}
\setlength{\parskip}{5pt plus 1pt}
\setlength{\emergencystretch}{3em}
\setcounter{tocdepth}{1}\setcounter{secnumdepth}{1}
\setlist[itemize]{leftmargin=3em,itemsep=4pt,topsep=4pt,parsep=0pt}
\allowdisplaybreaks
\newcommand{\N}{\mathbb{N}}
\newcommand{\Z}{\mathbb{Z}}
\newcommand{\Q}{\mathbb{Q}}
\newcommand{\R}{\mathbb{R}}
\newcommand{\C}{\mathbb{C}}
\newcommand{\F}{\mathbb{F}}
\newcommand{\A}{\mathbb{A}}
\newcommand{\PP}{\mathbb{P}}
\newcommand{\E}{\mathbb{E}}
\newcommand{\eps}{\varepsilon}
\newcommand{\dd}{\,\mathrm d}
\newcommand{\sref}[2]{\hyperlink{source:#1:#2}{[S#2]}}
\newif\ifshowcatalogue
\showcataloguetrue
\newcommand{\cataloguescope}[1]{\ifshowcatalogue
 \paragraph{Statement recorded in the supplied source.}
 #1\fi}
\begin{document}
% UnsolvedMath ID 2177; source catalogue code EP-560

\setcounter{section}{2176}

\section{The Size Ramsey Number of Balanced Complete Bipartite Graphs}\label{2177}

{\small\sffamily UnsolvedMath ID 2177 · Graph Theory}

\subsection{Definitions and mathematical statement}

\paragraph{Shared notation.}Unless a section says otherwise, $\N=\{1,2,\ldots\}$,
$\Z$ is the ring of integers, $\Q,\R,\C$ are the rational, real, and complex
fields, $[N]=\{1,\ldots,\lfloor N\rfloor\}$, and $\log$ is the natural
logarithm. For real functions of a parameter tending to infinity,
$f\ll g$ and $f=O(g)$ mean $|f|\leq Cg$ eventually for some fixed $C$;
$f\gg g$ reverses this comparison; $f=o(g)$ means $f/g\to0$;
$f\sim g$ means $f/g\to1$; and $f\asymp g$ means both order bounds.
A subscript on $O$ or $\ll$ specifies allowed constant dependence.
The expression $n^{o(1)}$ in an upper bound means growth smaller than
$n^\epsilon$ for each fixed $\epsilon>0$. Local definitions govern any
exception, finite-field convention, or use of the same symbol for another
object. An asymptotic claim is not automatically an effective algorithm.



The size Ramsey number minimizes host edge count, rather than host vertex count. The host is allowed to be any finite simple graph. In this entry the Ramsey property uses two edge colours and not necessarily induced copies.

A graph has a vertex set $V$ and edges that are unordered pairs of distinct vertices. A subgraph need not be induced unless that word is stated. A proper vertex colouring gives adjacent vertices different colours; $\chi(G)$ is the least number of colours.

$K_t$ is a complete graph and $C_t$ a simple cycle of length $t$. A complete multipartite graph has no edges within parts and all edges between different parts. An independent set has no internal edge; a clique has all its possible internal edges. \sref{2177}{1} \sref{2177}{2}

\paragraph{Statement recorded in the supplied source.}
Let $\hat{R}(G)$ denote the size Ramsey number, the minimal number of edges $m$ such that there is a graph $H$ with $m$ edges such that in any $2$-colouring of the edges of $H$ there is a monochromatic copy of $G$.
Determine $ \hat{R}(K_{n,n}), $ where $K_{n,n}$ is the complete bipartite graph with $n$ vertices in each component. \sref{2177}{1}

\paragraph{Residual task reported by the archive.}

The embedded review (2026-08-17) records: Determine the diagonal order of hat R(K\_{n,n}); the prominent conjecture is Theta(n\textasciicircum{}3 2\textasciicircum{}n). \sref{2177}{1}

\subsection{Short English statement}

What is the fewest host edges needed to force a monochromatic balanced complete bipartite graph under every two-colouring?

\subsection{Sources}

{\small

\begin{itemize}

\item[\hypertarget{source:2177:1}{S1}] ULAM.AI, \emph{UnsolvedMath}, supplied \texttt{problems.json}, record 2177 (EP-560), statement and background. The embedded literature-review date is 2026-08-17; that is the archive’s review date, not a fresh certification. Dataset landing page: \url{https://huggingface.co/datasets/ulamai/UnsolvedMath}.

\item[\hypertarget{source:2177:2}{S2}] T. F. Bloom, \emph{Erdős Problems}, Problem 560, \url{https://www.erdosproblems.com/560}. Reference imported from the supplied record; the live problem page was not independently checked for this compilation.

\item[\hypertarget{source:2177:3}{S3}] David Conlon, Jacob Fox, and Yuval Wigderson, Three early problems on size Ramsey numbers, Combinatorica 43 (2023), 743-768, \nolinkurl{doi:10.1007/s00493-023-00034-7.} \url{https://arxiv.org/abs/2111.05420}. Bibliographic information imported from the supplied record.

\item[\hypertarget{source:2177:4}{S4}] Erd\H{o}s, P. and Faudree, R. J. and Rousseau, C. C. and Schelp, R. H., The size Ramsey number. Period. Math. Hungar. (1978), 145-161. Bibliographic information imported from the supplied record.

\item[\hypertarget{source:2177:5}{S5}] Erd\H{o}s, P. and Rousseau, C. C., The size Ramsey number of a complete bipartite graph. Discrete Math. (1993), 259-262. Bibliographic information imported from the supplied record.

\item[\hypertarget{source:2177:6}{S6}] Ne\v{s}et\v{r}il, J. and R\"{o}dl, V., The structure of critical Ramsey graphs. Acta Math. Acad. Sci. Hungar. (1978), 295-300. Bibliographic information imported from the supplied record.

\end{itemize}

}

\label{end:2177}
\end{document}
